\section{Actual line-section transitions, finite \v{C}ech denominators and higher cohomology}
\label{sec:cech-line-native-cohomology}

Throughout, tensor powers, restrictions, pullbacks and \v{C}ech complexes
are the original native scheme-module constructions. For an invertible
$L$ and original $s:\mathcal O_X\to L$, write $V=D(s)$ for the complement
of the support of its actual zero ideal. The previously constructed
canonical frame on $V$ is used throughout. Write compositions in diagram
order. Definitions below record API constructions and carry no independent
source-theorem credit; the assertions about them are stated separately.

\begin{lemma}[Tensor multiplication in actual frames]
\label{lem:cech-line-native-two-frame-scalar}
\lean{EnvelopeProperty.SchemeModuleCechLineSectionLocalization.tensorSectionMap_tensorFrame}
\uses{lem:cech-line-section-original-tensor-values}
\leanok
Let $P,T$ be modules on $X$, let $e:P\simeq\mathcal O_X$ and
$d:T\simeq\mathcal O_X$ be actual frames, and let $t:\mathcal O_X\to T$.
Tensor multiplication by $t$, followed by the tensor frame induced by
$e,d$, is the whole map $e$ followed by multiplication by the actual
coefficient $\operatorname{frameScalar}(d,t)$.
\end{lemma}
\begin{proof}
The native tensor frame factors as $e\otimes\mathrm{id}_T$ followed by
the frame $d$ on the second factor and the unit constraint. Naturality
of tensor multiplication moves $e$ past the section map. The framed
section is its actual scalar endomorphism of $\mathcal O_X$; the unit
action on $\mathcal O_X$ is that same endomorphism. This proves equality
of entire module maps, without choosing or assuming a coefficient identity.
\end{proof}

\begin{definition}[The actual original twist system]
\label{def:cech-line-native-transitions}
\lean{EnvelopeProperty.SchemeModuleCechLineSectionLocalization.lineSectionStep,EnvelopeProperty.SchemeModuleCechLineSectionLocalization.twistStep,EnvelopeProperty.SchemeModuleCechLineSectionLocalization.twistIterate}
\uses{lem:cech-line-native-two-frame-scalar}
\leanok
For an original module $L$, section $s:\mathcal O_X\to L$ and module
$F$, put $F_m=F\otimes L^m$. The native tensor-power successor defines
$\ell_m:L^m\to L^{m+1}$ by multiplication by $s$, and
$T_m=\mathrm{id}_F\otimes\ell_m:F_m\to F_{m+1}$.
Define $T_{m,0}=\mathrm{id}_{F_m}$ and
$T_{m,a+1}=T_{m,a}T_{m+a}$. Thus $T_{m,a}$ has the literal target
$F\otimes L^{m+a}$ for every $a\geq0$.
These are constructions of the actual maps, rather than additional
vanishing or positivity theorems.
\end{definition}
\begin{proof}
Use the original section tensor map and the native isomorphism from the
monoidal tensor product to the tensor object in the defining recursion
for $L^{m+1}$. Tensor the resulting map by $F$. The recursion on $a$
then has the claimed source and target, including $a=0$.
\end{proof}

\begin{lemma}[Whole-map scalar formulas for transitions]
\label{lem:cech-line-native-transition-frame}
\lean{EnvelopeProperty.SchemeModuleCechLineSectionLocalization.twistStep_frame,EnvelopeProperty.SchemeModuleCechLineSectionLocalization.twistIterate_frame}
\uses{def:cech-line-native-transitions,lem:cech-line-native-two-frame-scalar}
\leanok
An actual frame $e:L\simeq\mathcal O_X$ identifies $T_m$ with
multiplication by $r=\operatorname{frameScalar}(e,s)$, between the
actual tensor frames of $F_m$ and $F_{m+1}$. It identifies $T_{m,a}$
with multiplication by $r^a$, between the actual frames of $F_m$ and
$F_{m+a}$. Both statements are equalities of whole original module maps.
\end{lemma}
\begin{proof}
Apply the two-frame tensor formula to the native power frame of $L^m$
and the original frame of $L$. Its recursive tensor comparison gives
the one-step formula. Induction on $a$ proves the iterate formula:
the identity corresponds to $r^0=1$, and composition of the actual
scalar actions multiplies $r^a$ by $r$.
\end{proof}

\begin{lemma}[Transitions in native pulled power frames]
\label{lem:cech-line-native-pulled-transition-frame}
\lean{EnvelopeProperty.SchemeModuleCechLineSectionLocalization.lineSectionStep_pullback_frame,EnvelopeProperty.SchemeModuleCechLineSectionLocalization.twistStep_pullback_frame}
\uses{lem:cech-line-native-two-frame-scalar,def:cech-line-native-transitions,lem:cech-line-section-original-native-pullback}
\leanok
For every morphism $g:Y\to X$ and actual frame
$e:g^*L\simeq\mathcal O_Y$, the pullback of $\ell_m$ has coefficient
$\operatorname{frameScalar}(e,g^*s)$ in the native pulled power frames.
After the genuine pullback tensor comparison, the pullback of $T_m$
has that same scalar action on $g^*F$.
\end{lemma}
\begin{proof}
Expand the native successor formula for the pulled power frame. The
pullback of the tensor-object isomorphism cancels its inverse.
Naturality of the actual section tensor map reduces the remaining
calculation to tensor multiplication in two genuine frames.
The monoidal pullback tensorator then gives the formula with $F$.
\end{proof}

\begin{lemma}[Every transition preserves canonical normalization]
\label{lem:cech-line-native-transition-normalized}
\lean{EnvelopeProperty.SchemeModuleCechLineSectionLocalization.twistStep_normalized,EnvelopeProperty.SchemeModuleCechLineSectionLocalization.twistIterate_normalized}
\uses{lem:cech-line-native-pulled-transition-frame,lem:cech-line-section-original-canonical-open}
\leanok
Suppose $L$ is invertible and write $V=D(s)$ for its literal nonzero
section open, with inclusion $\iota$. The canonical frame of $L|_V$
sends the original section to one. Consequently $\iota^*T_m$ and every
$\iota^*T_{m,a}$ preserve the actual canonical normalization
$\iota^*F_m\simeq\iota^*F$.
\end{lemma}
\begin{proof}
Specialize the pulled-frame transition formula to $g=\iota$ and the
canonical section frame. The coefficient is exactly $1$ by the actual
section-one law, so its scalar action is the identity. Induction on the
finite iterate, using functoriality of pullback, proves the second assertion.
\end{proof}

\begin{lemma}[Normalized values of original sections are unchanged]
\label{lem:cech-line-native-top-transition}
\lean{EnvelopeProperty.SchemeModuleCechLineSectionLocalization.normalizedTopValue_transition}
\uses{lem:cech-line-native-transition-normalized}
\leanok
For $y\in\Gamma(X,F_m)$ and every $a\geq0$, the canonical normalized
value on $D(s)$ of $T_{m,a}(y)$ equals that of $y$.
\end{lemma}
\begin{proof}
The actual pullback adjunction unit on global sections is natural in
the original module map $T_{m,a}$. Apply the whole-map normalization
identity to its image of $y$. This gives equality of the actual global
sections on $D(s)$.
\end{proof}

\begin{lemma}[Eventual affine lifts and zero clearing]
\label{lem:cech-line-native-affine-eventual}
\lean{EnvelopeProperty.SchemeModuleCechLineSectionLocalization.affine_eventual_normalized_lift,EnvelopeProperty.SchemeModuleCechLineSectionLocalization.affine_eventual_zero_clearing}
\uses{lem:cech-line-section-original-affine-normalized-lift,lem:cech-line-native-top-transition,lem:cech-line-native-transition-frame}
\leanok
Let $X$ be affine, let $F$ be quasicoherent, and give an ordinary actual
frame $e:L\simeq\mathcal O_X$. Every section of $F|_{D(s)}$ has a
canonical normalized lift in $\Gamma(X,F_m)$ for every sufficiently
large $m$. If $y\in\Gamma(X,F_m)$ has normalized value zero, then
$T_{m,a}(y)=0$ for every sufficiently large $a$.
The thresholds are produced internally from the original section and module.
\end{lemma}
\begin{proof}
The original affine normalized-lift theorem first gives a lift at one
exponent $N$. For every $m\geq N$, apply $T_{N,m-N}$ to that lift;
normalization preservation proves it is still a lift.
For zero clearing, the canonical normalization is an isomorphism, so
normalized zero implies ordinary restriction zero. The actual frame
identifies $D(s)$ with the basic open of its scalar $r$.
Quasicoherence of $F_m$ and genuine principal localization give some
$r^A y=0$. Every $r^a y$ with $a\geq A$ is then zero. The whole-map
transition-frame formula and the invertible frame identify this with
$T_{m,a}(y)=0$.
\end{proof}

\begin{definition}[The original section-open map]
\label{def:cech-line-native-open-map}
\lean{EnvelopeProperty.SchemeModuleCechLineSectionLocalization.sectionOpenMap}
\uses{lem:cech-line-section-original-affine-open}
\leanok
For $g:Y\to X$, set $V=D(s)$ and $W=D(g^*s)$. The actual section-open
map $h:W\to V$ is the unique open-immersion lift of $W\to Y\to X$
through $V\hookrightarrow X$. This is a canonical map construction.
\end{definition}
\begin{proof}
The previously proved equality $W=g^{-1}V$ puts the image of
$W\to Y\to X$ in the actual open $V$. Apply the native lift property
of that open immersion.
\end{proof}

\begin{lemma}[The actual section-open Cartesian square]
\label{lem:cech-line-native-open-square}
\lean{EnvelopeProperty.SchemeModuleCechLineSectionLocalization.sectionOpenMap_fac,EnvelopeProperty.SchemeModuleCechLineSectionLocalization.sectionOpenSquare}
\uses{def:cech-line-native-open-map,lem:cech-line-section-original-affine-open}
\leanok
The canonical map $h$ satisfies $h\iota_V=\iota_W g$, and the square
with sides $\iota_W,h,g,\iota_V$ is Cartesian. No square or open-matching
hypothesis is supplied.
\end{lemma}
\begin{proof}
The native lift construction gives the displayed triangle. The actual
inverse-image identity $D(g^*s)=g^{-1}D(s)$ identifies the open ranges;
the native open-immersion pullback criterion now proves Cartesianness.
\end{proof}

\begin{lemma}[Canonical power frames agree on the original square]
\label{lem:cech-line-native-open-power-frame}
\lean{EnvelopeProperty.SchemeModuleCechLineSectionLocalization.powerSectionFrame_openSquare}
\uses{lem:cech-line-native-open-square,lem:cech-line-section-original-canonical-open}
\leanok
For every $m$, the canonical frame of $(g^*L)^m$ on $W$ agrees with
the pullback of the canonical frame of $L^m$ on $V$, through the actual
pullback-power isomorphism and the canonical section-open square.
\end{lemma}
\begin{proof}
Pull back the actual original $s^m$. The native section-power pullback
formula identifies it with $(g^*s)^m$. Each proposed whole frame sends
this same section to one. This section is an isomorphism on $W$, so
cancellation proves equality of the entire frames.
\end{proof}

\begin{definition}[Whole original restriction and pushforward comparisons]
\label{def:cech-line-native-open-module-comparisons}
\lean{EnvelopeProperty.SchemeModuleCechLineSectionLocalization.sectionOpenModuleIso,EnvelopeProperty.SchemeModuleCechLineSectionLocalization.sectionOpenPushforwardIso}
\uses{lem:cech-line-native-open-square}
\leanok
Native pullback composition on the actual square defines
$h^*\iota_V^*F\simeq\iota_W^*g^*F$ for any original $F$.
If $g$ is an open immersion, the native open-square comparison further
defines $g^*\iota_{V*}\iota_V^*F\simeq\iota_{W*}\iota_W^*g^*F$.
These are specified whole-module isomorphisms, with no acyclicity assertion.
\end{definition}
\begin{proof}
Compose the native pullback-composition isomorphisms and the actual
triangle equality for the first construction. For the second, use the
existing open Cartesian pushforward comparison on the proved square,
then push forward the first whole-module isomorphism.
\end{proof}

\begin{lemma}[The actual open comparison preserves the adjunction unit]
\label{lem:cech-line-native-open-unit}
\lean{EnvelopeProperty.SchemeModuleCechLineSectionLocalization.sectionOpenPushforwardIso_unit}
\uses{def:cech-line-native-open-module-comparisons,lem:cech-line-native-open-square}
\leanok
For an open immersion $g$ and every original module $F$, pullback of
the unit $F\to\iota_{V*}\iota_V^*F$, followed by the actual open
pushforward comparison, equals the unit of $g^*F$ for $\iota_W$.
\end{lemma}
\begin{proof}
Identify the native open comparison with the mate of the actual
pullback square. The pullback-composition and section-open triangle
comparisons are used in that identity. The adjunction triangle law for
the mate gives the whole unit equality; the actual module comparison
and its inverse cancel. No sectionwise matching premise is used.
\end{proof}

\begin{definition}[Native twist comparison and normalized whole map]
\label{def:cech-line-native-normalized-pushforward}
\lean{EnvelopeProperty.SchemeModuleCechLineSectionLocalization.twistPullbackIso,EnvelopeProperty.SchemeModuleCechLineSectionLocalization.normalizedPushforwardMap}
\uses{lem:cech-line-section-original-canonical-open}
\leanok
Let $F_m=F\otimes L^m$. The native tensorator and pullback-power
isomorphism define $g^*F_m\simeq g^*F\otimes(g^*L)^m$.
Put $G=\iota_{V*}\iota_V^*F$. Define the actual normalized map
$\alpha_m:F_m\to G$ as the pullback-pushforward unit followed by the
pushforward of the canonical normalization $\iota_V^*F_m\simeq\iota_V^*F$.
These are native map constructions, not extra source-theorem credit.
\end{definition}
\begin{proof}
For the first map compose the inverse monoidal pullback tensorator with
the original pullback-power comparison on the second tensor factor.
For $\alpha_m$ use the native adjunction unit for the actual inclusion
$\iota_V$, and its pushforward functor on the already constructed
normalization map.
\end{proof}

\begin{lemma}[Actual normalized maps commute with original open pullback]
\label{lem:cech-line-native-normalization-pullback}
\lean{EnvelopeProperty.SchemeModuleCechLineSectionLocalization.powerSectionTensorMap_pullback,EnvelopeProperty.SchemeModuleCechLineSectionLocalization.normalizedTwistIso_openSquare,EnvelopeProperty.SchemeModuleCechLineSectionLocalization.normalizedPushforwardMap_pullback}
\uses{def:cech-line-native-normalized-pushforward,lem:cech-line-native-open-unit,lem:cech-line-native-open-power-frame,lem:cech-line-section-original-native-pullback}
\leanok
For every $g$, multiplication by $s^m$ pulls back, under the actual
twist comparison, to multiplication by $(g^*s)^m$. Canonical
normalization commutes with the actual section-open square as a whole
module map. If $g$ is an open immersion, $g^*\alpha_m$ followed by the
actual pushforward comparison equals the normalized map for $g^*F,g^*L,g^*s$
preceded by the actual twist comparison.
\end{lemma}
\begin{proof}
Native tensor-section naturality and the section-power pullback formula
prove the multiplication assertion. On $V$, multiplication by the
original $s^m$ is inverse to canonical normalization. Cancel that actual
isomorphism and use pullback-composition naturality to prove the
normalization-square equality. Finally expand $\alpha_m$: naturality
of the open pushforward comparison moves the normalization through it,
the preceding unit lemma handles the unit, and unit naturality handles
the actual twist comparison.
\end{proof}

\begin{lemma}[The normalized whole maps preserve every transition]
\label{lem:cech-line-native-pushforward-transition}
\lean{EnvelopeProperty.SchemeModuleCechLineSectionLocalization.normalizedPushforwardMap_transition}
\uses{def:cech-line-native-normalized-pushforward,lem:cech-line-native-transition-normalized}
\leanok
The actual maps satisfy $T_{m,a}\alpha_{m+a}=\alpha_m$ for every
$m,a\geq0$.
\end{lemma}
\begin{proof}
Expand each normalized map as its unit followed by the pushforward
normalization. Move $T_{m,a}$ past the unit by adjunction-unit
naturality. Inside the pushforward the actual transition-normalization
identity applies, giving the required equality of entire maps.
\end{proof}

\begin{lemma}[Actual eventual lifts over original affine opens]
\label{lem:cech-line-native-affine-open-lift}
\lean{EnvelopeProperty.SchemeModuleCechLineSectionLocalization.affineOpen_eventual_normalized_lift}
\uses{lem:cech-line-native-affine-eventual,lem:cech-line-native-normalization-pullback}
\leanok
Let $F$ be quasicoherent, let $U\subseteq X$ be an original affine
open, and give an actual frame of $L|_U$. Every
$x\in\Gamma(U,G)$ has, for every sufficiently large $m$, an actual
$y\in\Gamma(U,F_m)$ satisfying $\alpha_m(y)=x$.
\end{lemma}
\begin{proof}
Use the genuine open-section equivalence and the actual open
pushforward comparison to view $x$ as a section over $D(s|_U)$ of
the original pulled module. Apply the eventual affine lift theorem on
$U$. Transport the numerator back through the actual twist isomorphism
and open-section equivalence. The normalized-pullback whole-map
identity proves that this transported numerator has exactly value $x$.
\end{proof}

\begin{definition}[The native normalized and transition cochain maps]
\label{def:cech-line-native-cochain-maps}
\lean{EnvelopeProperty.SchemeModuleCechLineSectionLocalization.normalizedCechMap,EnvelopeProperty.SchemeModuleCechLineSectionLocalization.twistCechMap}
\uses{def:cech-line-native-transitions,def:cech-line-native-normalized-pushforward}
\leanok
For a family of original opens $\mathcal U=(U_i)$, applying the actual
native \v{C}ech complex functor to $\alpha_m$ gives
$A_m:C^\bullet(\mathcal U,F_m)\to C^\bullet(\mathcal U,G)$.
Applying it to $T_{m,a}$ gives
$B_{m,a}:C^\bullet(\mathcal U,F_m)\to C^\bullet(\mathcal U,F_{m+a})$.
These constructions retain the original native complexes and their differentials.
\end{definition}
\begin{proof}
Both maps are the functorial images of actual whole original module
maps. Hence they commute with the actual \v{C}ech differential. Their
component on each tuple intersection is evaluation of the respective
whole module map on that same intersection, by the native product projections.
\end{proof}

\begin{lemma}[A common actual twist lifts every finite native cochain]
\label{lem:cech-line-native-finite-cochain-lift}
\lean{EnvelopeProperty.SchemeModuleCechLineSectionLocalization.exists_finite_normalized_cochain_lift}
\uses{lem:cech-line-section-original-canonical-open,lem:cech-line-native-affine-open-lift,def:cech-line-native-cochain-maps}
\leanok
If $X$ is quasicompact and separated and $F$ is quasicoherent, there is
an internally chosen finite affine cover $\mathcal U$ with actual
frames of $L|_{U_i}$ such that for every degree $n$ and every
$x\in C^n(\mathcal U,G)$, the equation $A_m(y)=x$ has an actual
solution $y\in C^n(\mathcal U,F_m)$ for all sufficiently large $m$.
The bound may depend on $n$ and $x$.
\end{lemma}
\begin{proof}
Select the original finite framed affine cover. Separatedness makes
each actual tuple intersection affine; pull a frame from its first
factor. Apply the affine-open lifting theorem to each component of
$x$. There are finitely many tuples of length $n+1$, so the finite
maximum of their thresholds works for every component. The genuine
native \v{C}ech product isomorphism assembles their actual lifts.
Its projection formula for $A_m$ proves the whole cochain equality.
\end{proof}

\begin{lemma}[Finite transitions retain their actual pulled modules]
\label{lem:cech-line-native-transition-open-pullback}
\lean{EnvelopeProperty.SchemeModuleCechLineSectionLocalization.twistIterate_pullback_of_frame}
\uses{lem:cech-line-native-transition-frame,lem:cech-line-native-pulled-transition-frame,def:cech-line-native-normalized-pushforward}
\leanok
For every $g:Y\to X$ with an actual frame of $g^*L$, the native twist
isomorphisms conjugate $g^*T_{m,a}$ to the actual transition for
$g^*F,g^*L,g^*s$, with exactly exponents $m,m+a$.
\end{lemma}
\begin{proof}
The native pulled power frame factors through the actual power-pullback
isomorphism. The framed iterate formula on the source pullback and on
the pulled module has the same scalar $(g^*s)^a$. Monoidal pullback
naturality and cancellation of the genuine frames therefore identify
the entire finite transition maps.
\end{proof}

\begin{lemma}[The actual framed affine section-open preimage is affine]
\label{lem:cech-line-native-affine-section-preimage}
\lean{EnvelopeProperty.SchemeModuleCechLineSectionLocalization.sectionOpen_preimage_affine_of_frame}
\uses{lem:cech-line-section-original-affine-open}
\leanok
For an original affine open $U\subseteq X$ with an actual frame of
$L|_U$, the literal preimage $\iota_V^{-1}U$ is affine in $V=D(s)$.
\end{lemma}
\begin{proof}
On the original affine $U$, the pulled section-open is the basic open
of its actual coefficient, hence affine. Its image under the open
immersion $U\to X$ is exactly $U\cap V$. The native open-image and
preimage criterion then identifies $\iota_V^{-1}U$ as affine.
\end{proof}

\begin{lemma}[Zero clearing on original affine opens]
\label{lem:cech-line-native-affine-open-zero}
\lean{EnvelopeProperty.SchemeModuleCechLineSectionLocalization.affineOpen_eventual_zero_clearing}
\uses{lem:cech-line-native-affine-eventual,lem:cech-line-native-normalization-pullback,lem:cech-line-native-transition-open-pullback}
\leanok
Let $U$ be an original affine open with an actual frame of $L|_U$,
and let $F$ be quasicoherent. If $y\in\Gamma(U,F_m)$ satisfies
$\alpha_m(y)=0$, then $T_{m,a}(y)=0$ for every sufficiently large $a$.
\end{lemma}
\begin{proof}
The open-section equivalence and actual normalized-pullback identity
turn the hypothesis into normalized zero on the original affine $U$.
Apply genuine affine zero clearing there. The finite-transition
pullback formula identifies the resulting zero with the image of the
original $T_{m,a}(y)$. Both the twist comparison and the open-section
equivalence are injective, so the original section is zero.
\end{proof}

\begin{lemma}[A common actual transition clears a finite normalized-zero cochain]
\label{lem:cech-line-native-finite-cochain-zero}
\lean{EnvelopeProperty.SchemeModuleCechLineSectionLocalization.finite_normalized_cochain_zero}
\uses{lem:cech-line-native-affine-open-zero,def:cech-line-native-cochain-maps}
\leanok
Let $X$ be separated, $F$ quasicoherent, and $(U_i)$ a finite family
of affine opens with actual frames of $L|_{U_i}$. If
$y\in C^n(\mathcal U,F_m)$ has $A_m(y)=0$, then
$B_{m,a}(y)=0$ for every sufficiently large $a$.
No cover assumption is needed for this finite-family assertion.
\end{lemma}
\begin{proof}
The native product projection formula makes every tuple component of
$y$ normalized-zero on the same actual tuple intersection. The
intersection is affine and inherits a frame from its first factor.
Apply affine-open zero clearing componentwise and take the finite
maximum of the resulting bounds. The native cochain product
isomorphism detects zero, proving the asserted whole cochain equality.
\end{proof}

\begin{lemma}[Eventual original global cohomology gives actual native exactness]
\label{lem:cech-line-native-exact-from-global}
\lean{EnvelopeProperty.SchemeModuleCechLineSectionLocalization.exists_finite_normalized_cech_exactAt}
\uses{lem:cech-line-native-finite-cochain-lift,lem:cech-line-native-finite-cochain-zero,lem:cech-line-native-pushforward-transition,thm:scheme-module-cech-native-additive-comparison}
\leanok
Let $X$ be quasicompact and separated, $F$ quasicoherent and $q\geq0$.
Assume $H^{q+1}(X,F_m)=0$ for every sufficiently large $m$.
Then an internally chosen finite framed affine cover $\mathcal U$
has $C^\bullet(\mathcal U,G)$ exact in degree $q+1$.
\end{lemma}
\begin{proof}
Take a target cocycle $x$. Lift it to a cochain $y$ of an original
$F_m$, with $m$ above both the lifting bound and the global vanishing
bound. Since $A_m$ is a chain map, $dy$ has normalized value zero.
A sufficiently long actual transition kills $dy$, so
$y'=B_{m,a}(y)$ is an actual cocycle of $F_{m+a}$.
The global vanishing hypothesis applies at $m+a$. Native affine-cover
\v{C}ech comparison gives a genuine source boundary $y'=dz$.
Apply $A_{m+a}$; its chain-map identity and
$B_{m,a}A_{m+a}=A_m$ show $x=dA_{m+a}(z)$. The converse inclusion
is the actual equation $d^2=0$.
\end{proof}

\begin{theorem}[Actual higher vanishing on the original section open]
\label{thm:cech-line-native-section-open-vanishing}
\lean{EnvelopeProperty.SchemeModuleCechLineSectionLocalization.sectionOpen_restrict_H_succ_subsingleton}
\uses{lem:cech-line-native-exact-from-global,lem:cech-line-native-affine-section-preimage,lem:affine-pushforward-ordered-vanishing}
\leanok
Let $X$ be quasicompact and separated, $L$ invertible,
$s:\mathcal O_X\to L$ an original section, and $F$ quasicoherent.
For every $q\geq0$,
\[
 \bigl(\exists N\;\forall m\geq N,\ H^{q+1}(X,F\otimes L^m)=0\bigr)
 \quad\Longrightarrow\quad H^{q+1}(D(s),F|_{D(s)})=0.
\]
Here zero means that the actual native sheaf cohomology group is a
subsingleton. The finite cover, local denominators, comparisons and
source boundaries are constructed internally.
\end{theorem}
\begin{proof}
Use the preceding genuine exactness producer for the original
$G=\iota_{V*}\iota_V^*F$. The inverse images in $V$ of its selected
framed affine cover are affine and cover $V$. Native ordinary and
ordered \v{C}ech homology comparison, followed by the actual
pushforward-complex isomorphism for $\iota_V$, transfers the proved
exactness to this affine cover of the actual $\iota_V^*F$.
The native affine-cover comparison identifies its zero positive
homology with zero ordinary sheaf cohomology. Finally the whole native
restriction-pullback isomorphism identifies that module with the
literal original $F|_V$.
\end{proof}

The final theorem uses eventual vanishing on the original scheme itself.
No covering-space vanishing theorem alone supplies that hypothesis. In a
Nadel application, actual trace descent and the fixed original rounded
module comparison must first produce the stated original global groups.
